\chapter{Monte Carlo Simulation}
\label{ch:montecarlo}

\section{Concept}

Monte Carlo (MC) simulation quantifies how uncertainty in input parameters propagates through the model. In each iteration:
\begin{enumerate}
  \item Sample random values for selected parameters from their defined distributions
  \item Apply the sampled values to the MFA system
  \item Run the solver
  \item Record the results
\end{enumerate}

After $N$ iterations, the distribution of results shows the uncertainty range for all flows and stocks.

\section{Configuration (Sheet 4\_1)}

Uncertainty parameters are defined in the \texttt{4\_1\_Uncertainty\_Parameters} sheet:

\begin{table}[h]
\centering
\begin{tabular}{llllll}
\toprule
\textbf{Parameter\_Name} & \textbf{Distribution} & \textbf{Mean} & \textbf{StdDev} & \textbf{Min} & \textbf{Max} \\
\midrule
\texttt{TC\_05\_06} & normal & 0.65 & 0.05 & 0.4 & 0.9 \\
\texttt{F\_00\_01} & uniform & -- & -- & 900 & 1100 \\
\texttt{TC\_03\_04} & triangular & -- & -- & 0.3 & 0.5 \\
\bottomrule
\end{tabular}
\caption{MC parameter configuration. The \texttt{Mode} column (not shown) is required for triangular distributions.}
\end{table}

\subsection{Available Distributions}

\begin{description}
  \item[normal] Symmetric uncertainty. Parameters: \texttt{Mean}, \texttt{StdDev}. Optional: \texttt{Min}, \texttt{Max} for bounds.
  \item[uniform] Equal probability within range. Parameters: \texttt{Min}, \texttt{Max}.
  \item[triangular] Asymmetric with known most-likely value. Parameters: \texttt{Min}, \texttt{Max}, \texttt{Mode}.
  \item[lognormal] Always positive, right-skewed. Parameters: \texttt{Mean}, \texttt{StdDev}. Optional: \texttt{Min}, \texttt{Max}.
\end{description}

\section{Running MC Simulation}

In the notebook:

\begin{lstlisting}[style=python]
mc_results = run_monte_carlo(
    mfa_system_configured=mfa_system,
    uncertainty_params=uncertainty_params,
    n_iterations=100,
    process_logic_map=process_logic_map,
    flow_tc_map=flow_tc_map,
    dsm_params_config=dsm_params,
    fomp_params_config=fomp_params,
)
\end{lstlisting}

\begin{tipbox}
Start with 100 iterations to verify the setup, then increase to 500--1000 for publication-quality results. Each iteration runs the full solver, so computation time scales linearly.
\end{tipbox}

\section{TC Normalization}

When MC samples modify individual TCs, the outflow splits may no longer sum to 1.0. BioDYM uses \emph{rejection sampling}: if sampled TCs deviate too far from 100\%, the sample is rejected and redrawn. This preserves mass balance while respecting the defined uncertainty ranges.

\section{Results}

MC results include:
\begin{itemize}
  \item \textbf{Summary statistics} --- Mean, standard deviation, confidence intervals for all flows and stocks
  \item \textbf{Histograms} --- Distribution of key output variables across iterations
  \item \textbf{Sensitivity analysis} --- Tornado chart showing which parameters have the largest effect
  \item \textbf{Simulation paths} --- Spaghetti plots showing all iterations overlaid
  \item \textbf{Mass balance check} --- Per-process mass balance error across iterations (in scientific notation)
\end{itemize}
